% p10-01 例 1：矩形 ABCD AB=6 BC=8，P 在 AB 上从 A 出发，Q 在 BC 上从 B 出发
% 取某时刻 t=2: AP=2, BP=4, BQ=4 (△BPQ 等腰直角)
\documentclass[tikz,border=4pt]{standalone}
\usepackage{ctex}
\usepackage{amsmath}
\usetikzlibrary{calc, arrows.meta}
\begin{document}
\begin{tikzpicture}[scale=0.55, thick]
  \coordinate (A) at (0,0);
  \coordinate (B) at (6,0);
  \coordinate (C) at (6,8);
  \coordinate (D) at (0,8);
  \coordinate (P) at (2,0);   % AP=2
  \coordinate (Q) at (6,4);   % BQ=4

  \draw (A)--(B)--(C)--(D)--cycle;
  \draw[red, thick] (P)--(B)--(Q);
  \draw[red, thick] (P)--(Q);

  % 运动方向箭头（沿 AB 向 B；沿 BC 向 C）
  \draw[->, blue, thin] (1.5,0.5) -- (3.0,0.5);
  \draw[->, blue, thin] (5.5,3.5) -- (5.5,5.0);
  \node[blue, font=\footnotesize] at (2.2,0.95) {$P$ 向 $B$};
  \node[blue, font=\footnotesize] at (4.7,4.2) {$Q$ 向 $C$};

  \node[below left]  at (A) {$A$};
  \node[below right] at (B) {$B$};
  \node[above right] at (C) {$C$};
  \node[above left]  at (D) {$D$};
  \node[below] at (P) {$P$};
  \node[right] at (Q) {$Q$};

  \node[below] at ($(A)!0.5!(B)$) {\footnotesize $6$};
  \node[right] at ($(B)!0.5!(C)$) {\footnotesize $8$};
\end{tikzpicture}
\end{document}
